\documentclass[11pt]{article}
\usepackage{amsmath, amssymb}
\usepackage[margin=1in]{geometry}
\usepackage{hyperref}
\usepackage{booktabs}
\usepackage{listings}
\usepackage{xcolor}
\usepackage{graphicx}
\graphicspath{{figures/}}

\lstset{
  basicstyle=\ttfamily\small,
  breaklines=true,
  frame=single,
  backgroundcolor=\color{gray!8},
  showstringspaces=false,
  keywordstyle=\color{blue!60!black}\bfseries,
  commentstyle=\color{gray!60!black}\itshape,
  language=Python,
}

\title{Adaptive temporal subsampling for Zakharov rollouts\\
\large Event-aware per-case schedules for Tanaka collisions and Benjamin--Feir recurrence}
\author{}
\date{11 May 2026}

\begin{document}
\maketitle

\section*{Executive summary}

Production data is generated by integrating the Zakharov equations on a dense
time grid ($\Delta t = 0.08$, $t_{\max} = 200$ $\Rightarrow$ $T = 2501$ steps) and
saving $K = 200$ snapshots per case for FNO training. Uniform-$K$ subsampling
spreads those 200 snapshots evenly across the trajectory, but the events that
actually stress the surrogate --- soliton collisions in Tanaka, envelope
recurrence in Benjamin--Feir --- are short and locally complicated. With uniform
sampling those events fall between snapshots and are under-represented in
training. The fix is a per-case schedule that concentrates samples at events
while preserving full-trajectory coverage. Schedules are deterministic
(inverse-CDF over a quantile grid; no RNG), per-case (each trajectory picks its
own indices), and add negligible cost on top of the existing rollout.

\section{Motivation: why uniform fails}

\subsection{Conserved quantities are useless as event detectors}

The Zakharov Hamiltonian is conserved along trajectories, so any total-energy
diagnostic is flat in time and carries no event information. We need a
\emph{non-conserved} scalar that spikes during the focusing / collision /
recurrence events we care about.

\subsection{FNO state sensitivity}

Empirically, FNO rollouts diverge fastest at the moments where the underlying
PDE state changes rapidly --- soliton overlap, BF envelope peaks. These regions
are exactly where the training mix is thinnest under uniform-$K$ sampling: a
Tanaka collision is $\mathcal{O}(1)$ wide on a 200\,s trajectory, so uniform
sampling lands maybe 1 sample on it; BF recurrence cycles span $\sim 50$\,s and
again get under-sampled at $K = 200$.

\section{Activity signal}

\subsection{Two regime-specific choices}

We use two different non-conserved signals depending on the regime:

\begin{itemize}
  \item \textbf{Tanaka (soliton collisions).} Surface-gradient energy
  \[
    S(t) = \|\eta_x(t)\|^2 = \int_0^L \eta_x(x,t)^2 \, dx,
  \]
  computed via spectral derivative. This dips during destructive tail overlap
  before a collision (since the inner product
  $\int \eta_{1,x}\,\eta_{2,x}\,dx$ becomes negative as solitons of opposite
  orientation overlap) and spikes at peak focusing. The dip--spike pattern
  means $|dS/dt|$ is large on \emph{both} sides of the collision, which is
  what we want for biasing samples toward the event.

  \item \textbf{Benjamin--Feir (envelope recurrence).} Envelope peak
  \[
    M(t) = \|\eta(t)\|_{\infty} = \max_x |\eta(x,t)|.
  \]
  Surface-gradient energy is dominated by the fast carrier and washes out the
  slow envelope modulation; the envelope peak tracks recurrence directly, with
  $M(t)$ humping at each Akhmediev focusing.
\end{itemize}

Both signals are scalar per (case, time), computed on-device from the
$(T, B, N)$ eta trajectory just after the rollout finishes.

\subsection{Memory-aware computation}

The naive spectral derivative for $S(t)$ takes a full $(T, B, N)$ complex
FFT intermediate, which at $T = 2501$, $B = 256$, $N = 1024$ in float64 is
$\sim 20$\,GB and OOMs on a 48\,GB card alongside the eta/xi buffers. We compute
it instead with \texttt{jax.lax.scan} over the time axis so the FFT intermediate
is only $(B, N)$ per step:

\begin{lstlisting}
@jax.jit
def grad_energy_traj(eta_traj, length):
    n = eta_traj.shape[-1]
    k = 2.0 * jnp.pi * jnp.fft.fftfreq(n, d=length / n)
    def step(_, eta_t):
        eta_hat = jnp.fft.fft(eta_t, axis=-1)
        eta_x = jnp.real(jnp.fft.ifft(1j * k * eta_hat, axis=-1))
        return None, jnp.sum(eta_x ** 2, axis=-1)
    _, s_TB = jax.lax.scan(step, None, eta_traj)
    return jnp.swapaxes(s_TB, 0, 1)  # (B, T)
\end{lstlisting}

The envelope-peak signal is just \texttt{jnp.max(jnp.abs(...), axis=-1)} and
needs no scan.

\section{Density construction}

For each case $b$ we build a sampling density over the dense time grid
$\{t_0, \ldots, t_{T-1}\}$:
\[
  \mathrm{density}_b(t) \;=\; (1 - \alpha)\,\mathrm{uniform}(t) \;+\; \alpha \, \mathrm{smooth}\bigl(\mathrm{activity}_b(t)\bigr),
\]
where $\alpha \in [0, 1]$ mixes a uniform prior (always covers the full
trajectory) against an activity-weighted bias. The smoothing is Gaussian with
width $\sigma$ (in dense-time steps) to avoid spiky over-concentration on a
single step.

\paragraph{Two ways to build activity.}
We provide two modes for turning a raw signal into a density weight:

\begin{itemize}
  \item \texttt{density\_mode = "rate"}: $\mathrm{activity}(t) = |dS/dt| / (S + \varepsilon)$.
  Picks up localized state changes; correct for Tanaka collisions, which are
  short bursts on top of a quiet baseline.

  \item \texttt{density\_mode = "magnitude"}: $\mathrm{activity}(t) = S(t) - \min_t S(t)$.
  Picks up where the signal is high; correct for BF, where the envelope
  modulates continuously and we want to bias toward whichever times have the
  largest peaks rather than the moments of fastest change.
\end{itemize}

\paragraph{Power sharpening.}
For shallow modulation (BF envelope typically only moves $\sim 20\%$ from
trough to peak), the magnitude density is nearly uniform on its own. We
sharpen by raising it to a power:
\[
  \mathrm{activity}(t) \mapsto \bigl(\mathrm{activity}(t)\bigr)^p,
\]
applied after subtracting the per-case minimum (so the baseline doesn't
dominate after exponentiation). For the BF regen $p = 4$ gives clear
three-cluster concentration on the envelope peaks; $p = 1$ stays flat.

\section{Deterministic inverse-CDF sampling}

Given the per-case density, we pick $K$ indices via inverse-CDF at $K$
equispaced quantiles:
\begin{align}
  \mathrm{cdf}(t) &= \mathrm{cumsum}(\mathrm{density} / \textstyle\sum \mathrm{density}), \\
  q_k &= (k + \tfrac{1}{2}) / K, \quad k = 0, \ldots, K-1, \\
  \mathrm{idx}_k &= \mathrm{searchsorted}(\mathrm{cdf}, q_k).
\end{align}
We then pin endpoints ($\mathrm{idx}_0 = 0$, $\mathrm{idx}_{K-1} = T-1$) and
sweep once to enforce strict monotonicity in case adjacent quantiles land on
the same index. No RNG: same trajectory always produces the same indices, and
two cases with identical activity profiles get identical schedules.

The result is an $(B, K)$ \texttt{int32} array of per-case indices into the
dense time grid.

\section{Plumbing into data generation}

\subsection{End-to-end per-batch flow}

For each batch of $B = 256$ cases:

\begin{enumerate}
  \item \textbf{Rollout.} \texttt{batched\_rollout} on the dense grid produces
  \texttt{eta\_TBN} and \texttt{xi\_TBN} of shape $(T, B, N)$, kept on-device.

  \item \textbf{Activity signal.} One of \texttt{grad\_energy\_traj} (scanned
  over time) or \texttt{envelope\_peak\_traj} reduces eta to a $(B, T)$ scalar
  field. Transferred to host.

  \item \textbf{Index selection.} \texttt{adaptive\_indices\_from\_signal} on
  host builds the per-case density, integrates the CDF, and picks
  $\mathrm{idx} \in \mathbb{Z}^{B \times K}$.

  \item \textbf{Gather.} \texttt{jax.numpy.take\_along\_axis} slices eta and xi
  at the per-case indices, producing $(B, K, N)$ tensors.

  \item \textbf{Sparse gxi.} $G(\eta)\xi$ is computed via spectral DNO series
  only at the $K$ selected times (not the dense $T$), in chunks to keep memory
  bounded. This is a major win: the dense-grid gxi pass would cost $T / K
  \approx 12\times$ more.

  \item \textbf{Flatten and write.} The $(B, K, N)$ tensors flatten to
  $(BK, N)$ for the npz, alongside the $(B, K)$ index array and the dense-time
  vector so each sample's actual time is recoverable.
\end{enumerate}

\subsection{Why per-case schedules}

A single shared schedule (the same $K$ time indices for every case in a batch)
would be uniform from each case's point of view, because different cases peak
at different times: a 2-soliton collision happens at the case-specific time
when the solitons overlap, BF recurrence period scales with $1/\varepsilon^2$
per case. Per-case scheduling adds nothing to the runtime budget --- one tiny
\texttt{searchsorted} per case after the rollout --- and is the whole point.

\section{Parameter choices for the regen runs}

\begin{table}[h]
\centering
\begin{tabular}{lcc}
\toprule
                  & Tanaka                      & BF                          \\
\midrule
signal            & $\|\eta_x\|^2$ (grad\_energy) & $\|\eta\|_{\infty}$ (envelope) \\
density mode      & rate                        & magnitude                   \\
$\alpha$          & 0.5                         & 0.85                        \\
$\sigma$ (steps)  & 50                          & 20                          \\
power             & 1.0                         & 4.0                         \\
\bottomrule
\end{tabular}
\caption{Adaptive sampling parameters used in the production regeneration of
\texttt{tanaka\_2\_adaptive\_*} and \texttt{bf\_2\_adaptive\_*}.}
\end{table}

\paragraph{Why $\alpha$ differs.} $\alpha = 0.5$ on Tanaka leaves substantial
uniform coverage because much of the trajectory between collisions is still
worth sampling (the FNO sees plenty of clean propagation, but collisions are
where it fails). $\alpha = 0.85$ on BF concentrates more aggressively on
recurrence peaks because the envelope modulation \emph{is} the only structure
to lock onto --- there's no quiet baseline to subsidize uniformly.

\paragraph{Why $\sigma$ differs.} $\sigma = 50$ on Tanaka smooths the
$|dS/dt|/S$ rate signal over $\sim 4$\,s of dense-grid time, broadening the
collision peak so we sample its approach and aftermath, not just the apex.
$\sigma = 20$ on BF is enough because the envelope humps are themselves slow
features that don't need much smoothing.

\section{Smoke-test results}

The adaptive sampler is implemented in
\texttt{solver/gen\_data/adaptive\_sampling.py} and exercised by two demos:

\begin{itemize}
  \item \texttt{solver/evals/adaptive\_subsample\_demo.py} --- 2-crest Tanaka
  at $h = 0.05$. Renders the $\eta(x, t)$ heatmap with both uniform and
  adaptive sample times overlaid, plus the actual saved $(\eta, \xi)$ at each
  adaptive snapshot. Confirms that the rate density clusters samples around
  the collision event and leaves uniform coverage of the propagating regime.

  \item \texttt{solver/evals/adaptive\_subsample\_demo\_bf.py} --- canonical
  BF setup ($n_{\text{carr}} = 10$, $\varepsilon_c = 0.12$, side\_offset = 3,
  $h = 1.5$, $t_{\max} = 150$). Confirms that the magnitude density with
  $\alpha = 0.85$, $p = 4$ produces three sample clusters aligned with the
  three envelope-peak recurrence events visible in $\|\eta\|_{\infty}(t)$,
  whereas $p = 1$ produces nearly uniform samples because the modulation only
  spans $\sim 20\%$.
\end{itemize}

\section{Generator CLI}

Both \texttt{generate\_tanaka\_dataset\_v2} and \texttt{generate\_bf\_dataset}
accept the following flags:

\begin{lstlisting}[language=bash]
--adaptive                          # opt in
--adaptive_signal {envelope, grad_energy}
--adaptive_density_mode {rate, magnitude}
--adaptive_alpha     <float in [0, 1]>
--adaptive_smooth_sigma <float>     # Gaussian sigma in dense-time steps
--adaptive_power     <float>        # 1.0 = none; raise for sharper bias
\end{lstlisting}

Without \texttt{--adaptive}, the generator falls back to the existing uniform
stride logic. With \texttt{--adaptive}, the per-batch index array is also
written to the npz under \texttt{subsample\_indices\_\{batch\_tag\}.npy} so each
saved sample's source dense-grid time is fully recoverable post-hoc.

\section{Cost}

The adaptive machinery adds, per batch:

\begin{itemize}
  \item One \texttt{scan} reduction of eta over $T$ steps (grad\_energy path)
  or one \texttt{max} reduction (envelope path) --- both well under 1\,s on
  A100 / RTX 6000 Ada at $T = 2501$, $B = 256$, $N = 1024$.

  \item One host-side density+CDF+searchsorted per case --- microseconds.

  \item Gather + sparse gxi pass replacing the dense gxi pass --- this is a
  \emph{savings}, not a cost, since the dense pass evaluated $G(\eta)\xi$ at
  all $T$ time slices.
\end{itemize}

Net effect on wall time per batch: zero to mildly negative.

\end{document}